
\documentclass{article}
\usepackage{colortbl}
\usepackage{makecell}
\usepackage{multirow}
\usepackage{supertabular}

\begin{document}

\newcounter{utterance}

\centering \large Interaction Transcript for game `adventuregame', experiment `home\_deliver\_three\_basic\_hard\_invlimittwo', episode 4 with reka{-}flash{-}3.1{-}t0.0.
\vspace{24pt}

{ \footnotesize  \setcounter{utterance}{1}
\setlength{\tabcolsep}{0pt}
\begin{supertabular}{c@{$\;$}|p{.15\linewidth}@{}p{.15\linewidth}p{.15\linewidth}p{.15\linewidth}p{.15\linewidth}p{.15\linewidth}}
    \# & \multicolumn{2}{c}{Player} && \multicolumn{2}{c}{Game Master} \\
    \hline

    \theutterance \stepcounter{utterance}  
    & & & \multicolumn{4}{p{0.6\linewidth}}{
        \cellcolor[rgb]{0.9,0.9,0.9}{
            \makecell[{{p{\linewidth}}}]{
                \texttt{\tiny{[P1$\langle$GM]}}
                \texttt{You are playing a text adventure game. I will describe what you can perceive in the game. You write the single action you want to take in the game starting with >. Only reply with an action.} \\
\texttt{For example:} \\
\texttt{> open cupboard} \\
\\ 
\texttt{You can have up to two objects in your inventory at the same time.} \\
\\ 
\texttt{Your goal for this game is: Put the pillow in the refrigerator, the book in the wardrobe and the plate on the shelf.} \\
\\ 
\texttt{Once you have achieved your goal, write "> done" to end the game.} \\
\\ 
\texttt{You are in a living room now. There are a table, a couch, a chair and a side table. There is a passage to a hallway here.} \\
            }
        }
    }
     \\ \\

    \theutterance \stepcounter{utterance}  
    & \multicolumn{4}{p{0.6\linewidth}}{
        \cellcolor[rgb]{0.561,0.737,0.561}{
            \makecell[{{p{\linewidth}}}]{
                \texttt{\tiny{[P1$\rangle$GM]}}
                \texttt{<reasoning>} \\
\texttt{Let's see. The goal is to put the pillow in the refrigerator, the book in the wardrobe, and the plate on the shelf. Right now, I'm in the living room with a table, couch, chair, side table, and a passage to the hallway. The inventory can hold up to two objects. So first, I need to figure out where each item is located.} \\
\\ 
\texttt{Wait, the problem statement mentions the pillow, book, and plate, but the initial description doesn't list them. Maybe they are on the table, side table, or somewhere else. Wait, the user hasn't provided the initial items. Wait, maybe the initial description is incomplete. Wait, the user's message says "You are in a living room now. There are a table, a couch, a chair and a side table. There is a passage to a hallway here." But where are the pillow, book, and plate? Are they mentioned here? Let me check again.} \\
\\ 
\texttt{No, the initial description only lists the furniture. So perhaps the items are in the environment but not explicitly stated. Wait, maybe the user expects me to assume that the items are present in the room. Or maybe the items are in the hallway or other areas. Wait, the problem says "You can perceive in the game" but the user hasn't provided the current items. Wait, maybe the user is going to describe the environment step by step, and I need to act based on that. Wait, the initial message says "You are playing a text adventure game. I will describe what you can perceive in the game. You write the single action you want to take in the game starting with >. Only reply with an action."} \\
\\ 
\texttt{So the user will describe the current room's contents, and I have to act based on that. But in the initial message, the user wrote: "You are in a living room now. There are a table, a couch, a chair and a side table. There is a passage to a hallway here." So the items (pillow, book, plate) must be in this room. Wait, but the initial description doesn't mention them. That's confusing. Maybe the user made a mistake, or perhaps the items are in the room but not listed. Wait, maybe the user expects me to ask for clarification, but the instructions say to only reply with an action. Hmm.} \\
\\ 
\texttt{Wait, perhaps the items are present but not explicitly mentioned. For} \\
            }
        }
    }
    & & \\ \\

    \theutterance \stepcounter{utterance}  
    & & & \multicolumn{2}{p{0.3\linewidth}}{
        \cellcolor[rgb]{0.9,0.9,0.9}{
            \makecell[{{p{\linewidth}}}]{
                \texttt{\tiny{[GM$|$GM]}}
                \texttt{Round: 0} \\
            }
        }
    }
    & & \\ \\

    \theutterance \stepcounter{utterance}  
    & & & \multicolumn{2}{p{0.3\linewidth}}{
        \cellcolor[rgb]{0.9,0.9,0.9}{
            \makecell[{{p{\linewidth}}}]{
                \texttt{\tiny{[GM$|$GM]}}
                \texttt{<reasoning>} \\
\texttt{Let's see. The goal is to put the pillow in the refrigerator, the book in the wardrobe, and the plate on the shelf. Right now, I'm in the living room with a table, couch, chair, side table, and a passage to the hallway. The inventory can hold up to two objects. So first, I need to figure out where each item is located.} \\
\\ 
\texttt{Wait, the problem statement mentions the pillow, book, and plate, but the initial description doesn't list them. Maybe they are on the table, side table, or somewhere else. Wait, the user hasn't provided the initial items. Wait, maybe the initial description is incomplete. Wait, the user's message says "You are in a living room now. There are a table, a couch, a chair and a side table. There is a passage to a hallway here." But where are the pillow, book, and plate? Are they mentioned here? Let me check again.} \\
\\ 
\texttt{No, the initial description only lists the furniture. So perhaps the items are in the environment but not explicitly stated. Wait, maybe the user expects me to assume that the items are present in the room. Or maybe the items are in the hallway or other areas. Wait, the problem says "You can perceive in the game" but the user hasn't provided the current items. Wait, maybe the user is going to describe the environment step by step, and I need to act based on that. Wait, the initial message says "You are playing a text adventure game. I will describe what you can perceive in the game. You write the single action you want to take in the game starting with >. Only reply with an action."} \\
\\ 
\texttt{So the user will describe the current room's contents, and I have to act based on that. But in the initial message, the user wrote: "You are in a living room now. There are a table, a couch, a chair and a side table. There is a passage to a hallway here." So the items (pillow, book, plate) must be in this room. Wait, but the initial description doesn't mention them. That's confusing. Maybe the user made a mistake, or perhaps the items are in the room but not listed. Wait, maybe the user expects me to ask for clarification, but the instructions say to only reply with an action. Hmm.} \\
\\ 
\texttt{Wait, perhaps the items are present but not explicitly mentioned. For} \\
            }
        }
    }
    & & \\ \\

    \theutterance \stepcounter{utterance}  
    & & & \multicolumn{2}{p{0.3\linewidth}}{
        \cellcolor[rgb]{0.9,0.9,0.9}{
            \makecell[{{p{\linewidth}}}]{
                \texttt{\tiny{[GM$|$GM]}}
                \texttt{command\_tag\_missing} \\
            }
        }
    }
    & & \\ \\

    \theutterance \stepcounter{utterance}  
    & & & \multicolumn{2}{p{0.3\linewidth}}{
        \cellcolor[rgb]{0.9,0.9,0.9}{
            \makecell[{{p{\linewidth}}}]{
                \texttt{\tiny{[GM$|$GM]}}
                {'goal_states_achieved': [], 'game_successfully_finished': False}
            }
        }
    }
    & & \\ \\

    \theutterance \stepcounter{utterance}  
    & & & \multicolumn{2}{p{0.3\linewidth}}{
        \cellcolor[rgb]{0.9,0.9,0.9}{
            \makecell[{{p{\linewidth}}}]{
                \texttt{\tiny{[GM$|$GM]}}
                \texttt{command\_tag\_missing} \\
            }
        }
    }
    & & \\ \\

\end{supertabular}
}

\end{document}
